% Source: https://artofproblemsolving.com/community/c3h1605750_proof_that_1__2__3__4___is_not_112
% Stable identity: aops-topic-1605750-q01
\textbf{(a)} The conclusion is true, but the argument is not a proof.
\begin{itemize}
\item $\infty$ is not a real number, so it cannot be substituted into the formula $\dfrac{n(n+1)}{2}$, which is proved only for positive integers $n$. Expressions such as $\infty^2+\infty+\dfrac{1}{6}$ and ``the roots of a quadratic in $\infty$'' have no meaning in the real numbers, so the algebra done with them proves nothing.
\item The argument also never says what the ``sum'' of an infinite series is. Without a definition, neither the claim $1+2+3+\cdots=-\dfrac{1}{12}$ nor its denial can be proved or disproved.
\end{itemize}

\textbf{(b)} By definition, a series converges to a real number $L$ if its partial sums $s_n$ get arbitrarily close to $L$: for every $\varepsilon>0$ there is an $N$ with $|s_n-L|<\varepsilon$ for all $n\ge N$. Here
\[
s_n=1+2+\cdots+n=\dfrac{n(n+1)}{2}\ge\dfrac{n^2}{2}.
\]
Suppose the series converged to some real $L$. Taking $\varepsilon=1$, all $s_n$ with $n\ge N$ would satisfy $s_n<L+1$. But $s_n\ge\dfrac{n^2}{2}>L+1$ as soon as $n>\sqrt{2|L|+2}$, a contradiction. So the series converges to no real number; its partial sums grow without bound (it diverges to $+\infty$). In particular it does not converge to $-\dfrac{1}{12}$; even more simply, every partial sum is at least $1$, so they could never approach a negative number. $\blacksquare$

(The value $-\dfrac{1}{12}$ arises in advanced mathematics from methods that assign values to some divergent series, such as the analytic continuation of the zeta function, $\zeta(-1)=-\dfrac{1}{12}$. It is not the sum of the series in the ordinary sense.)
